% Generated draft: requires mathematical and source review.
\documentclass[11pt]{article}
\usepackage[margin=25mm]{geometry}
\usepackage{fontspec}
\setmainfont{Latin Modern Roman}
\usepackage{amsmath,amssymb,mathtools,unicode-math}
\setmathfont{Latin Modern Math}
\usepackage{enumitem,needspace,xurl,hyperref}
\setlength{\parindent}{0pt}
\setlength{\parskip}{5pt}
\setlength{\emergencystretch}{3em}
\newcommand{\N}{\mathbb{N}}
\newcommand{\Z}{\mathbb{Z}}
\newcommand{\Q}{\mathbb{Q}}
\newcommand{\R}{\mathbb{R}}
\newcommand{\C}{\mathbb{C}}
\newcommand{\F}{\mathbb{F}}
\newcommand{\A}{\mathbb{A}}
\newcommand{\PP}{\mathbb{P}}
\newcommand{\E}{\mathbb{E}}
\newcommand{\eps}{\varepsilon}
\newcommand{\dd}{\,\mathrm{d}}
\begin{document}
\section*{UnsolvedMath ID: 30004265}
\textbf{Source problem code:} catalogue entry
\subsection{Definitions and mathematical statement}
The supplied record does not specify the quiver, ground field, representation, dimension vector, group action, or intersection-cohomology conventions. Consequently, it determines a comparison problem but does not determine a single precise conjecture.

For terminology only, suppose that \(Q=(Q_0,Q_1)\) is a finite quiver, \(k\) is a field, and \(M=((M_i)_{i\in Q_0},(M_a)_{a\in Q_1})\) is a finite-dimensional representation of \(Q\) over \(k\). Thus an arrow \(a:i\to j\) has a linear map \(M_a:M_i\to M_j\). For a dimension vector \(\mathbf e=(e_i)_{i\in Q_0}\) with \(0\leq e_i\leq\dim_k M_i\), the quiver Grassmannian \(\operatorname{Gr}_{\mathbf e}(M)\) parametrizes subrepresentations of dimension vector \(\mathbf e\); its \(k\)-points are
\[
\operatorname{Gr}_{\mathbf e}(M)(k)=\left\{(U_i)_{i\in Q_0}: U_i\subseteq M_i,\ \dim_k U_i=e_i,\ M_a(U_i)\subseteq U_j\text{ for every }a:i\to j\right\}.
\]
These are explanatory parameters, not hypotheses supplied for the intended problem. Given an algebraic group action on such a variety, an orbit \(\mathcal O\) has a Zariski closure \(\overline{\mathcal O}\). The record denotes by \(f_{\mathcal O}\) an intersection-cohomology function associated with this closure, but supplies no definition of its values, coefficient theory, normalization, or extension outside the closure.

\textbf{Source comparison problem.} Determine how the family \(\{f_{\mathcal O}\}\), described in the source record as a canonical basis arising from quiver-Grassmannian orbit closures, is related to variants of Lusztig's canonical basis in quantized enveloping algebras.

The intended orbit family, vector space or algebra having \(\{f_{\mathcal O}\}\) as a basis, quantized enveloping algebra, canonical-basis variants, and comparison maps are unspecified. In particular, the record does not specify whether the desired relationship is an identification, an image under a map, or another compatibility. No particular equality, correspondence, or universally quantified assertion can therefore be extracted without additional source information. The description of the family as a canonical basis is retained as source terminology, rather than asserted here as an established basis theorem.
\subsection{Short English statement}
The problem asks how intersection-cohomology functions on quiver-Grassmannian orbit closures compare with variants of Lusztig's canonical basis. The supplied record leaves the geometric data, function conventions, algebra, and comparison unspecified, so a precise conjecture requires further source clarification. Its open status is supplied dated curation.
\subsection{Sources}
\begin{itemize}
\item Local registry: \texttt{data/unsolvedmath/problems.json}, permanent UnsolvedMath ID \(30004265\), title \emph{Canonical Bases from Quiver Grassmannian Orbit Closures}; exact source problem code \texttt{catalogue entry}. The supplied collection is \emph{Oberwolfach Reports Open Problems}. No individual proposer is identified.
\item Supplied report citation: \emph{Mini-Workshop: Degeneration Techniques in Representation Theory} (2020), \emph{Oberwolfach Reports}, pp.\ 2869--2909, DOI \texttt{10.4171/owr/2019/46}. \url{https://doi.org/10.4171/owr/2019/46}. The record separately gives proposed year 2019; both year indications are retained without resolving their relationship.
\item Supplied literature source, labeled \emph{OWR 2019 46}: \url{https://ems.press/content/serial-article-files/46825}.
\item Literature assessment dated exactly \texttt{2026-08-22}: the supplied curation labels the problem open and says that the relationship is explicitly posed as an open problem, with no general identification located. This is supplied dated curation, not an independently verified literature assessment. The background requests independent verification of time-sensitive status and cited claims; no source visit or such verification is claimed here.
\end{itemize}
\end{document}
